\clearpage
\VerdierPhysicalPageAnchor{0198}%
\begingroup\small\setlength{\parskip}{0pt}\setlength{\emergencystretch}{1.5em}
\setlength{\abovedisplayskip}{4pt plus 2pt minus 2pt}%
\setlength{\belowdisplayskip}{4pt plus 2pt minus 2pt}%

\VerdierSemanticAnchor{ch3d-2-3-1-proof-b-conclusion}
et par suite $v$ est un épimorphisme, ce qui démontre la propriété $(E_3)$
et achève la démonstration de la proposition.

\VerdierSemanticAnchor{ch3d-2-3-2}
\medskip\noindent\textbf{Proposition 2.3.2.}
\textit{Soit $M$ un ensemble d'objets de $\mathcal A$ possédant les
propriétés $(E_1)$, $(E_2)$ et $(E_3)$. Soit
$Y\in\operatorname{Ob}\operatorname{comp}(\mathcal A)$ un complexe tel que,
pour tout entier $n$, $H^n(Y)\in M$ et tel qu'il existe un entier $n_0$ tel
que $H^n(Y)=0$, pour $n\geq n_0$. Il existe alors un complexe
$X\in\operatorname{Ob}\operatorname{comp}^{-}(\mathcal A)$ dont les
composants sont des éléments de $M$ et un quasi-isomorphisme :}
\[
\dot u\colon X\longrightarrow Y.
\]

\VerdierSemanticAnchor{ch3d-2-3-2-proof}
Tout d'abord, il existe (1.2.1) un complexe
$Y'\in\operatorname{Ob}\operatorname{comp}^{-}(\mathcal A)$ et un
quasi-isomorphisme $Y'\longrightarrow Y$. On peut donc se ramener au cas où
$Y$ est un objet de $\operatorname{comp}^{-}(\mathcal A)$. Nous allons
construire par récurrence sur $n$, une suite de morphismes :

\VerdierSemanticAnchor{ch3d-2-3-2-induction}
\[
u_n\colon X_n\longrightarrow X_{n-1}
\]
possédant les propriétés suivantes :

\VerdierSemanticAnchor{ch3d-2-3-2-induction-properties}
\noindent 1) Pour $n$ assez petit $X_n=Y$ et $u_n=\operatorname{id}$.

\noindent 2) Les composants de degré $i\geq -n$ du complexe $X_n$ sont des
éléments de $M$. Les noyaux et images des différentielles de degré
$i\geq -n$ de $X_n$ sont des éléments de $M^\wedge$ (2.1.1).

\noindent 3) Le morphisme $u_n$ est un quasi-isomorphisme. Il est isomorphe
au morphisme $p((\dot u_n)^{-n},-n)$ (2.1.1.3).

\VerdierSemanticAnchor{ch3d-2-3-2-limits}
Ceci fait, il suffit de poser $X=\varprojlim_n X_n$ et de prendre pour
morphisme $\dot u$ le morphisme canonique :
\[
\dot u\colon X=\varprojlim_n X_n\longrightarrow\varinjlim_n X_n=Y.
\]

\VerdierSemanticAnchor{ch3d-2-3-2-induction-step}
Pour $n$ assez petit, on pose $X_n=Y$ et $\dot u_n=\operatorname{id}_Y$.
Supposons que l'on ait construit pour tout entier $n<q$ des complexes $X_n$
et des morphismes $\dot u_n$ possédant les propriétés 1), 2) et 3) ci-dessus. Le
complexe $X_{q-1}$ est de la forme :

\VerdierSemanticAnchor{ch3d-2-3-2-xqminus1}
\[
X_{q-1}=\cdots\longrightarrow Y^{-q-2}\longrightarrow Y^{-q-1}
\longrightarrow Z^{-q}\longrightarrow L^{-q+1}\longrightarrow L^{-q+2}
\longrightarrow\cdots,
\]
où les $Y^i$ sont les composants du complexe $Y$ et les $L^i$ sont des
éléments de $M$. On a alors un épimorphisme :

\VerdierSemanticAnchor{ch3d-2-3-2-quotient}
\[
Z^{-q}\longrightarrow
Z^{-q}/\operatorname{Im}(Y^{-q-1}\longrightarrow Z^{-q})
\longrightarrow0
\]
\endgroup
